\documentclass[12pt]{article}
\usepackage[T1]{fontenc}
\usepackage[utf8]{inputenc}
\usepackage{lmodern}
\usepackage{textcomp}
\usepackage{enumitem}
\usepackage{geometry}
\geometry{margin=1in}
\begin{document}
\begin{center}
Answer Key - Philosophy - Graduate Level Assessment 18 \
\small (Answers are ordered exactly as in the question paper.)
\end{center}

\section*{Multiple Choice}
\begin{enumerate}[label=\arabic*.]
\item \textbf{Answer:} A
\\ \textit{Options:} A) do that which is good and not to approve of it.; B) be motivated by genuine benevolence.; C) approve of the actions of others without understanding them.; D) act without a sense of self-interest.; E) feel empathy without taking action.
\\ \textit{Note:} Butler argues that moral goodness inherently involves approval; therefore, one cannot genuinely perform a good act without also recognizing and endorsing its goodness, highlighting the inseparability of moral action and moral judgment.
\item \textbf{Answer:} A
\\ \textit{Options:} A) affirms the antecedent or denies the consequent; B) leads to an invalid conclusion; C) generalizes a specific case; D) includes three or more alternatives; E) contradicts the conclusion
\\ \textit{Note:} A valid disjunctive syllogism requires a major premise that asserts one of the disjuncts is true or that the other disjunct must be false, thus affirming the antecedent or denying the consequent to reach a logical conclusion.
\item \textbf{Answer:} A
\\ \textit{Options:} A) 333 BCE; B) 1066 CE; C) 535 BCE; D) 515 BCE; E) 1914 CE
\\ \textit{Note:} The Babylonian captivity, also known as the Exile, began in 586 BCE, not 333 BCE, as it was during this time that the Babylonians conquered Jerusalem and exiled many Judeans, making 586 BCE the pivotal date in this historical and philosophical context.
\item \textbf{Answer:} E
\\ \textit{Options:} A) Roman Catholic; B) Lutherans; C) Methodists; D) Greek Orthodox; E) Quakers
\\ \textit{Note:} The Quakers, or the Religious Society of Friends, advocate for a simplified and unadorned form of Christianity that emphasizes personal experience of the divine, inner spirituality, and direct communion with God, devoid of elaborate rituals and hierarchies.
\item \textbf{Answer:} B
\\ \textit{Options:} A) vice is in our power, but virtue is not.; B) virtue is in our power, and so is vice.; C) vice is in our power, and so is virtue.; D) both virtue and vice are not in our power.; E) virtue is in our power, but vice is not.
\\ \textit{Note:} Aristotle asserts that both virtue and vice are a result of our choices and actions, emphasizing that moral character is within our control and dependent on our decisions, thus making the statement correct.
\end{enumerate}

\section*{True / False}
\begin{enumerate}[label=\arabic*.]
\setcounter{enumi}{5}
\item \textbf{Answer:} False
\\ \textit{Options:} A) True; B) False
\\ \textit{Note:} The statement is false because the moral status of an individual is not universally determined by their actions or convictions, and many philosophers argue that moral degradation is a complex issue influenced by various factors, including societal, psychological, and contextual elements.
\item \textbf{Answer:} True
\\ \textit{Options:} A) True; B) False
\\ \textit{Note:} SCNT accurately refers to somatic cell nuclear transfer, which is a scientific method employed in cloning that involves transferring the nucleus of a somatic cell into an enucleated egg cell to create an organism with the same genetic material as the donor.
\item \textbf{Answer:} True
\\ \textit{Options:} A) True; B) False
\\ \textit{Note:} Lukianoff and Haidt contend that trigger warnings reflect a mindset that shifts responsibility for emotional distress from the individual to external sources, thereby fostering a culture of victimhood and avoidance rather than resilience and personal accountability.
\item \textbf{Answer:} True
\\ \textit{Options:} A) True; B) False
\\ \textit{Note:} Nussbaum argues that the identification of specific virtues can lead to conflicts between diverse cultural and philosophical perspectives, as differing societies may prioritize different virtues based on their unique contexts and values.
\item \textbf{Answer:} True
\\ \textit{Options:} A) True; B) False
\\ \textit{Note:} Ross critiques utilitarianism for prioritizing outcomes over the moral intentions behind actions, arguing that this oversight undermines the complexity of ethical decision-making.
\end{enumerate}

\section*{Long Form Response}
\begin{enumerate}[label=\arabic*.]
\setcounter{enumi}{10}
\item \textbf{Answer:} Velleman argues that the option of euthanasia can impose a profound psychological burden on patients, compelling them to justify their continued existence. This burden arises from the societal and personal pressures to demonstrate one's worthiness of life, particularly in the face of suffering or diminished quality of life. Patients may feel that their lives require validation through productivity or happiness, leading to feelings of inadequacy or despair if they cannot meet these expectations. Consequently, rather than providing relief, the option of euthanasia may exacerbate feelings of alienation and worthlessness, ultimately undermining the intrinsic value of human life. In this way, Velleman highlights the ethical implications of euthanasia, suggesting that it may inadvertently shift the focus from the sanctity of life to a utilitarian calculus of existence.
\\ \textit{Note:} correct because it captures Velleman's argument that the availability of euthanasia can generate an ethical dilemma for patients, who may feel compelled to justify their existence based on societal standards of worth, leading to psychological distress and a potential sense of inadequacy.
\item \textbf{Answer:} The straw person fallacy occurs when an individual misrepresents an opponent's argument by attributing to them a position that is either exaggerated or entirely fabricated, which is then easily refuted. This tactic allows the arguer to claim victory over a distorted version of the opponent's stance, rather than engaging with the actual argument presented. The implications of this fallacy for logical discourse are significant, as it undermines the integrity of debate by shifting focus away from substantive issues and fostering misunderstanding. Consequently, it hinders productive dialogue, erodes trust between interlocutors, and diminishes the overall quality of argumentative exchange in philosophical discourse. Engaging with genuine arguments is essential for advancing knowledge and understanding, making the straw person fallacy a detrimental practice in rigorous argumentation.
\\ \textit{Note:} The answer correctly identifies the straw person fallacy as a misrepresentation of an opponent's argument, emphasizing that it allows for an easier refutation of distorted claims rather than addressing the actual position, which ultimately hinders honest and constructive discourse.
\end{enumerate}

\vfill
\noindent\textit{End of Answer Key}
\end{document}